\chapter{svnlook Reference---Subversion Repository
Examination}\label{svn.ref.svnlook}

\texttt{svnlook} is a command-line utility for examining different
aspects of a Subversion repository. It does not make any changes to the
repository---it\textquotesingle s just used for ``peeking''.
\texttt{svnlook} is typically used by the repository hooks, but a
repository administrator might find it useful for diagnostic purposes.

Since \texttt{svnlook} works via direct repository access (and thus can
be used only on the machine that holds the repository), it refers to the
repository with a path, not a URL.

If no revision or transaction is specified, \texttt{svnlook} defaults to
the youngest (most recent) revision of the repository.

Options in \texttt{svnlook} are global, just as they are in \texttt{svn}
and \texttt{svnadmin}; however, most options apply to only one
subcommand since the functionality of \texttt{svnlook} is
(intentionally) limited in scope:

\protect\phantomsection\label{svn.ref.svnlook.sw}{}

\begin{description}
\item[\protect\phantomsection\label{svn.ref.svnlook.sw.copy_info}{}\texttt{-\/-copy-info}]
Causes \texttt{svnlook\ changed} to show detailed copy source
information.
\item[\protect\phantomsection\label{svn.ref.svnlook.sw.diff_cmd}{}\texttt{-\/-diff-cmd}
\textless CMD\textgreater{}]
Specifies an external program to use to show differences between files.
When \texttt{svnlook\ diff} is invoked without this option, it uses
Subversion\textquotesingle s internal differencing engine, which
provides unified diffs by default. If you want to use an external
differencing program, use \texttt{-\/-diff-cmd}. You can then pass
options to the specified program using the \texttt{-\/-extensions}
(\texttt{-x}) option.
\item[\protect\phantomsection\label{svn.ref.svnlook.sw.diff_copy_from}{}\texttt{-\/-diff-copy-from}]
Print differences for copied items against the copy source.
\item[\protect\phantomsection\label{svn.ref.svnlook.sw.extensions}{}\texttt{-\/-extensions}
(\texttt{-x}) \textless ARG\textgreater{}]
Specifies customizations which Subversion should make when performing
difference calculations. Valid extensions include:

\begin{description}
\item[\texttt{-\/-ignore-space-change} (\texttt{-b})]
Ignore changes in the amount of white space.
\item[\texttt{-\/-ignore-all-space} (\texttt{-w})]
Ignore all white space.
\item[\texttt{-\/-ignore-eol-style}]
Ignore changes in EOL (end-of-line) style.
\item[\texttt{-\/-show-c-function} (\texttt{-p})]
Show C function names in the diff output.
\item[\texttt{-\/-unified} (\texttt{-u})]
Show three lines of unified diff context.
\end{description}

The default value is \texttt{-u}.

Note that when Subversion is configured to invoke an external diff
command, the value of the \texttt{-\/-extensions} (\texttt{-x}) option
isn\textquotesingle t restricted to the previously mentioned options,
but may be \emph{any} additional arguments which Subversion should pass
to that command. If you wish to pass multiple arguments, you must
enclose all of them in quotes.
\item[\protect\phantomsection\label{svn.ref.svnlook.sw.full_paths}{}\texttt{-\/-full-paths}]
Causes \texttt{svnlook\ tree} to display full paths instead of
hierarchical, indented path components.
\item[\protect\phantomsection\label{svn.ref.svnlook.sw.ignore_properties}{}\texttt{-\/-ignore-properties}]
Instructs \texttt{svnlook\ diff} to suppress output of property changes.
\item[\protect\phantomsection\label{svn.ref.svnlook.sw.limit}{}\texttt{-\/-limit}
(\texttt{-l}) \textless ARG\textgreater{}]
Limit output to a maximum number of \textless ARG\textgreater{} items.
\item[\protect\phantomsection\label{svn.ref.svnlook.sw.no_diff_deleted}{}\texttt{-\/-no-diff-deleted}]
Prevents \texttt{svnlook\ diff} from printing differences for deleted
files. The default behavior when a file is deleted in a
transaction/revision is to print the same differences that you would see
if you had left the file but removed all the content.
\item[\protect\phantomsection\label{svn.ref.svnlook.sw.no_diff_added}{}\texttt{-\/-no-diff-added}]
Prevents \texttt{svnlook\ diff} from printing differences for added
files. The default behavior when you add a file is to print the same
differences that you would see if you had added the entire contents of
an existing (empty) file.
\item[\protect\phantomsection\label{svn.ref.svnlook.sw.non_recursive}{}\texttt{-\/-non-recursive}
(\texttt{-N})]
Operate on a single directory only.
\item[\protect\phantomsection\label{svn.ref.svnlook.sw.properties_only}{}\texttt{-\/-properties-only}]
Instructs \texttt{svnlook\ diff} to show only property changes.
\item[\protect\phantomsection\label{svn.ref.svnlook.sw.revision}{}\texttt{-\/-revision}
(\texttt{-r}) \textless REV\textgreater{}]
Specifies a particular revision number that you wish to examine.
\item[\protect\phantomsection\label{svn.ref.svnlook.sw.revprop}{}\texttt{-\/-revprop}]
Operates on a revision property instead of a property specific to a file
or directory. This option requires that you also pass a revision with
the \texttt{-\/-revision} (\texttt{-r}) option.
\item[\protect\phantomsection\label{svn.ref.svnlook.sw.show_inherited_props}{}\texttt{-\/-show-inherited-props}]
Works with \texttt{svnlook\ propget} and \texttt{svnlook\ proplist} to
display the versioned properties inherited by a path.
\item[\protect\phantomsection\label{svn.ref.svnlook.sw.transaction}{}\texttt{-\/-transaction}
(\texttt{-t}) \textless ID\textgreater{}]
Specifies a particular transaction ID that you wish to examine.
\item[\protect\phantomsection\label{svn.ref.svnlook.sw.show_ids}{}\texttt{-\/-show-ids}]
Shows the filesystem node revision IDs for each path in the filesystem
tree.
\item[\protect\phantomsection\label{svn.ref.svnlook.sw.verbose}{}\texttt{-\/-verbose}
(\texttt{-v})]
Be verbose. When used with \texttt{svnlook\ proplist}, for example, this
causes Subversion to display not just the list of properties, but their
values also.
\item[\protect\phantomsection\label{svn.ref.svnlook.sw.xml}{}\texttt{-\/-xml}]
Prints output in XML format.
\end{description}

\section{svnlook author}\label{svn.ref.svnlook.c.author}

\emph{Print the author.}

\textbf{Synopsis}

\texttt{svnlook\ author\ REPOS\_PATH}

\subsection{Description}

Print the author of a revision or transaction in the repository.

\subsection{Options}

\begin{verbatim}
--revision (-r) REV
--transaction (-t) ID
\end{verbatim}

\subsection{Examples}

\texttt{svnlook\ author} is handy, but not very exciting:

\begin{verbatim}
$ svnlook author -r 40 /var/svn/repos 
sally
\end{verbatim}

\section{svnlook cat}\label{svn.ref.svnlook.c.cat}

\emph{Print the contents of a file.}

\textbf{Synopsis}

\texttt{svnlook\ cat\ REPOS\_PATH\ PATH\_IN\_REPOS}

\subsection{Description}

Print the contents of a file.

\subsection{Options}

\begin{verbatim}
--revision (-r) REV
--transaction (-t) ID
\end{verbatim}

\subsection{Examples}

This shows the contents of a file in transaction \texttt{ax8}, located
at \texttt{/trunk/README}:

\begin{verbatim}
$ svnlook cat -t ax8 /var/svn/repos /trunk/README

               Subversion, a version control system.
               =====================================

$LastChangedDate: 2003-07-17 10:45:25 -0500 (Thu, 17 Jul 2003) $

Contents:

     I. A FEW POINTERS
    II. DOCUMENTATION
   III. PARTICIPATING IN THE SUBVERSION COMMUNITY
…
\end{verbatim}

\section{svnlook changed}\label{svn.ref.svnlook.c.changed}

\emph{Print the paths that were changed.}

\textbf{Synopsis}

\texttt{svnlook\ changed\ REPOS\_PATH}

\subsection{Description}

Print the paths that were changed in a particular revision or
transaction, as well as ``svn update-style'' status letters in the first
two columns:

\begin{description}
\item[\texttt{\textquotesingle{}A\ \textquotesingle{}}]
Item added to repository
\item[\texttt{\textquotesingle{}D\ \textquotesingle{}}]
Item deleted from repository
\item[\texttt{\textquotesingle{}U\ \textquotesingle{}}]
File contents changed
\item[\texttt{\textquotesingle{}\_U\textquotesingle{}}]
Properties of item changed; note the leading underscore
\item[\texttt{\textquotesingle{}UU\textquotesingle{}}]
File contents and properties changed
\end{description}

Files and directories can be distinguished, as directory paths are
displayed with a trailing ``\texttt{/}'' character.

\subsection{Options}

\begin{verbatim}
--copy-info
--revision (-r) REV
--transaction (-t) ID
\end{verbatim}

\subsection{Examples}

This shows a list of all the changed files and directories in revision
39 of a test repository. Note that the first changed item is a
directory, as evidenced by the trailing \texttt{/}:

\begin{verbatim}
$ svnlook changed -r 39 /var/svn/repos
A   trunk/vendors/deli/
A   trunk/vendors/deli/chips.txt
A   trunk/vendors/deli/sandwich.txt
A   trunk/vendors/deli/pickle.txt
U   trunk/vendors/baker/bagel.txt
_U  trunk/vendors/baker/croissant.txt
UU  trunk/vendors/baker/pretzel.txt
D   trunk/vendors/baker/baguette.txt
\end{verbatim}

Here\textquotesingle s an example that shows a revision in which a file
was renamed:

\begin{verbatim}
$ svnlook changed -r 64 /var/svn/repos
A   trunk/vendors/baker/toast.txt
D   trunk/vendors/baker/bread.txt
\end{verbatim}

Unfortunately, nothing in the preceding output reveals the connection
between the deleted and added files. Use the \texttt{-\/-copy-info}
option to make this relationship more apparent:

\begin{verbatim}
$ svnlook changed -r 64 --copy-info /var/svn/repos
A + trunk/vendors/baker/toast.txt
    (from trunk/vendors/baker/bread.txt:r63)
D   trunk/vendors/baker/bread.txt
\end{verbatim}

\section{svnlook date}\label{svn.ref.svnlook.c.date}

\emph{Print the datestamp.}

\textbf{Synopsis}

\texttt{svnlook\ date\ REPOS\_PATH}

\subsection{Description}

Print the datestamp of a revision or transaction in a repository.

\subsection{Options}

\begin{verbatim}
--revision (-r) REV
--transaction (-t) ID
\end{verbatim}

\subsection{Examples}

This shows the date of revision 40 of a test repository:

\begin{verbatim}
$ svnlook date -r 40 /var/svn/repos/
2003-02-22 17:44:49 -0600 (Sat, 22 Feb 2003)
\end{verbatim}

\section{svnlook diff}\label{svn.ref.svnlook.c.diff}

\emph{Print differences of changed files and properties.}

\textbf{Synopsis}

\texttt{svnlook\ diff\ REPOS\_PATH}

\subsection{Description}

Print GNU-style differences of changed files and properties in a
repository.

\subsection{Options}

\begin{verbatim}
--diff-cmd CMD
--diff-copy-from
--ignore-properties
--no-diff-added
--no-diff-deleted
--properties-only
--revision (-r) REV
--transaction (-t) ID
--extensions (-x) ARG
\end{verbatim}

\subsection{Examples}

This shows a newly added (empty) file, a modified binary file, and a
renamed (that is, copied and deleted) file with modifications:

\begin{verbatim}
$ svnlook diff -r 40 /var/svn/repos
Copied: trunk/relish.txt (from rev 39, trunk/vendors/deli/pickle.txt)
===================================================================
--- trunk/relish.txt                            (rev 0)
+++ trunk/relish.txt    2013-01-29 20:39:17 UTC (rev 40)
@@ -0,0 +1 @@
+Pickle relish is mostly made from cucumbers.

Deleted: trunk/vendors/deli/pickle.txt
===================================================================
--- trunk/vendors/deli/pickle.txt                           (rev 39)
+++ trunk/vendors/deli/pickle.txt   2013-01-29 20:39:17 UTC (rev 49)
@@ -1 +0,0 @@
-Pickles are mostly made from cucumbers.

Modified: trunk/vendors/deli/logo.jpg
===================================================================
(Binary files differ)

Added: trunk/vendors/deli/soda.txt
===================================================================
$
\end{verbatim}

By default, \texttt{svnlook\ diff} will treat copied files very much
like any other added file, displaying in their entirety the contents of
the new file and merely using a different label to draw the copy/add
distinction. However, you can use the \texttt{-\/-diff-copy-from} option
to cause \texttt{svnlook\ diff} to consider a copied file as worthy of
mention only if it differs from the file from which it was copied, and
to actually describe those differences.

\begin{verbatim}
$ svnlook diff -r 40 /var/svn/repos --diff-copy-from
Copied: trunk/relish.txt (from rev 39, trunk/vendors/deli/pickle.txt)
===================================================================
--- trunk/vendors/deli/pickle.txt   2013-01-29 20:39:17 UTC (rev 39)
+++ trunk/relish.txt    2013-01-29 20:47:40 UTC (rev 3)
@@ -1 +1 @@
-Pickles are mostly made from cucumbers.
+Pickle relish is mostly made from cucumbers.

Deleted: trunk/vendors/deli/pickle.txt
===================================================================
--- trunk/vendors/deli/pickle.txt   (rev 39)
+++ trunk/vendors/deli/pickle.txt   2013-01-29 20:39:17 UTC (rev 40)
@@ -1 +0,0 @@
-Pickles are mostly made from cucumbers.

Modified: trunk/vendors/deli/logo.jpg
===================================================================
(Binary files differ)

Added: trunk/vendors/deli/soda.txt
==============================================================================
$
\end{verbatim}

Use the \texttt{-\/-no-diff-deleted} option to silence output regarding
deleted files.

\begin{verbatim}
$ svnlook diff -r 40 /var/svn/repos --no-diff-deleted
Copied: trunk/relish.txt (from rev 39, trunk/vendors/deli/pickle.txt)
===================================================================
--- trunk/relish.txt                            (rev 0)
+++ trunk/relish.txt    2013-01-29 20:39:17 UTC (rev 40)
@@ -0,0 +1 @@
+Pickle relish is mostly made from cucumbers.

Modified: trunk/vendors/deli/logo.jpg
===================================================================
(Binary files differ)

Added: trunk/vendors/deli/soda.txt
==============================================================================
$
\end{verbatim}

Note that in each of the previous examples, when a file has a nontextual
\texttt{svn:mime-type} property, the differences are not explicitly
shown.

\section{svnlook dirs-changed}\label{svn.ref.svnlook.c.dirs-changed}

\emph{Print the directories that were themselves changed.}

\textbf{Synopsis}

\texttt{svnlook\ dirs-changed\ REPOS\_PATH}

\subsection{Description}

Print the directories that were themselves changed (property edits) or
whose file children were changed.

\subsection{Options}

\begin{verbatim}
--revision (-r) REV
--transaction (-t) ID
\end{verbatim}

\subsection{Examples}

This shows the directories that changed in revision 40 in our sample
repository:

\begin{verbatim}
$ svnlook dirs-changed -r 40 /var/svn/repos
trunk/vendors/deli/
\end{verbatim}

\section{svnlook filesize}\label{svn.ref.svnlook.c.filesize}

\emph{Print the size (in bytes) of a versioned file.}

\textbf{Synopsis}

\texttt{svnlook\ filesize\ REPOS\_PATH\ PATH\_IN\_REPOS}

\subsection{Description}

Print the file size (in bytes) of the file located at
\textless PATH\_IN\_REPOS\textgreater{} in the \texttt{HEAD} revision of
the repository identified by \textless REPOS\_PATH\textgreater{} as a
base-10 integer followed by an end-of-line character. Use the
\texttt{-\/-revision} (\texttt{-r}) and \texttt{-\/-transaction}
(\texttt{-t}) options to specify a version of the file other than
\texttt{HEAD} whose file size you wish to display.

\subsection{Options}

\begin{verbatim}
--revision (-r) REV
--transaction (-t) ID
\end{verbatim}

\subsection{Examples}

The following demonstrates how to display the size of the
\texttt{trunk/vendors/deli/soda.txt} file as it appeared in revision 40
of our sample repository:

\begin{verbatim}
$ svnlook filesize -r 40 /var/svn/repos trunk/vendors/deli/soda.txt
23
$
\end{verbatim}

\section{svnlook help (h, ?)}\label{svn.ref.svnlook.c.help}

\emph{Help!}

\textbf{Synopsis}

\texttt{Also\ svnlook\ -h\ and\ svnlook\ -?.}

Also \texttt{svnlook\ -h} and \texttt{svnlook\ -?}.

\subsection{Description}

Displays the help message for \texttt{svnlook}. This command, like its
brother, \texttt{svn\ help}, is also your friend, even though you never
call it anymore and forgot to invite it to your last party.

\subsection{Options}

None

\section{svnlook history}\label{svn.ref.svnlook.c.history}

\emph{Print information about the history of a path in the repository
(or the root directory if no path is supplied).}

\textbf{Synopsis}

\texttt{svnlook\ history\ REPOS\_PATH\ {[}PATH\_IN\_REPOS{]}}

\subsection{Description}

Print information about the history of a path in the repository (or the
root directory if no path is supplied).

\subsection{Options}

\begin{verbatim}
--limit (-l) ARG
--revision (-r) REV
--show-ids
\end{verbatim}

\subsection{Examples}

This shows the history output for the path \texttt{/branches/bookstore}
as of revision 13 in our sample repository:

\begin{verbatim}
$ svnlook history -r 13 /var/svn/repos /branches/bookstore --show-ids
REVISION   PATH <ID>
--------   ---------
      13   /branches/bookstore <1.1.r13/390>
      12   /branches/bookstore <1.1.r12/413>
      11   /branches/bookstore <1.1.r11/0>
       9   /trunk <1.0.r9/551>
       8   /trunk <1.0.r8/131357096>
       7   /trunk <1.0.r7/294>
       6   /trunk <1.0.r6/353>
       5   /trunk <1.0.r5/349>
       4   /trunk <1.0.r4/332>
       3   /trunk <1.0.r3/335>
       2   /trunk <1.0.r2/295>
       1   /trunk <1.0.r1/532>
\end{verbatim}

\section{svnlook info}\label{svn.ref.svnlook.c.info}

\emph{Print the author, datestamp, log message size, and log message.}

\textbf{Synopsis}

\texttt{svnlook\ info\ REPOS\_PATH}

\subsection{Description}

Print the author, datestamp, log message size (in bytes), and log
message, followed by a newline character.

\subsection{Options}

\begin{verbatim}
--revision (-r) REV
--transaction (-t) ID
\end{verbatim}

\subsection{Examples}

This shows the info output for revision 40 in our sample repository:

\begin{verbatim}
$ svnlook info -r 40 /var/svn/repos
sally
2003-02-22 17:44:49 -0600 (Sat, 22 Feb 2003)
16
Rearrange lunch.
\end{verbatim}

\section{svnlook lock}\label{svn.ref.svnlook.c.lock}

\emph{If a lock exists on a path in the repository, describe it.}

\textbf{Synopsis}

\texttt{svnlook\ lock\ REPOS\_PATH\ PATH\_IN\_REPOS}

\subsection{Description}

Print all information available for the lock at
\textless PATH\_IN\_REPOS\textgreater. If
\textless PATH\_IN\_REPOS\textgreater{} is not locked, print nothing.

\subsection{Options}

None

\subsection{Examples}

This describes the lock on the file \texttt{tree.jpg}:

\begin{verbatim}
$ svnlook lock /var/svn/repos tree.jpg
UUID Token: opaquelocktoken:ab00ddf0-6afb-0310-9cd0-dda813329753
Owner: harry
Created: 2005-07-08 17:27:36 -0500 (Fri, 08 Jul 2005)
Expires: 
Comment (1 line):
Rework the uppermost branches on the bald cypress in the foreground.
\end{verbatim}

\section{svnlook log}\label{svn.ref.svnlook.c.log}

\emph{Print the log message, followed by a newline character.}

\textbf{Synopsis}

\texttt{svnlook\ log\ REPOS\_PATH}

\subsection{Description}

Print the log message.

\subsection{Options}

\begin{verbatim}
--revision (-r) REV
--transaction (-t) ID
\end{verbatim}

\subsection{Examples}

This shows the log output for revision 40 in our sample repository:

\begin{verbatim}
$ svnlook log -r 40 /var/svn/repos/
Rearrange lunch.
\end{verbatim}

\section{svnlook propget (pget, pg)}\label{svn.ref.svnlook.c.propget}

\emph{Print the raw value of a property on a path in the repository.}

\textbf{Synopsis}

\texttt{svnlook\ propget\ REPOS\_PATH\ PROPNAME\ {[}PATH\_IN\_REPOS{]}}

\subsection{Description}

List the value of a property on a path in the repository.

\subsection{Options}

\begin{verbatim}
--revision (-r) REV
--revprop
--show-inherited-props
--transaction (-t) ID
\end{verbatim}

\subsection{Examples}

This shows the value of the ``seasonings'' property on the file
\texttt{/trunk/sandwich} in the \texttt{HEAD} revision:

\begin{verbatim}
$ svnlook pg /var/svn/repos seasonings /trunk/sandwich
mustard
\end{verbatim}

This shows the inherited values of the \texttt{svn:auto-props} property
on the directory \texttt{/trunk} in revision 14:

\begin{verbatim}
$ svnlook pg repos svn:auto-props trunk --show-inherited-props -v -r14
Inherited properties on '/trunk',
from '/':
  svn:auto-props
    *.py = svn:eol-style=native
    *.c = svn:eol-style=native
    *.h = svn:eol-style=native
\end{verbatim}

\section{svnlook proplist (plist, pl)}\label{svn.ref.svnlook.c.proplist}

\emph{Print the names and values of versioned file and directory
properties.}

\textbf{Synopsis}

\texttt{svnlook\ proplist\ REPOS\_PATH\ {[}PATH\_IN\_REPOS{]}}

\subsection{Description}

List the properties of a path in the repository. With
\texttt{-\/-verbose} (\texttt{-v}), show the property values too.

\subsection{Options}

\begin{verbatim}
--revision (-r) REV
--revprop
--show-inherited-props
--transaction (-t) ID
--verbose (-v)
--xml
\end{verbatim}

\subsection{Examples}

This shows the names of properties set on the file
\texttt{/trunk/README} in the \texttt{HEAD} revision:

\begin{verbatim}
$ svnlook proplist /var/svn/repos /trunk/README
  original-author
  svn:mime-type
\end{verbatim}

This is the same command as in the preceding example, but this time
showing the property values as well:

\begin{verbatim}
$ svnlook -v proplist /var/svn/repos /trunk/README
  original-author : harry
  svn:mime-type : text/plain
\end{verbatim}

This shows the properties inherited by a directory:

\begin{verbatim}
$ svnlook pl /var/svn/repos branches --show-inherited-props -v
Inherited properties on '/branches',
from '/':
  svn:auto-props
    *.py = svn:eol-style=native
    *.c = svn:eol-style=native
    *.h = svn:eol-style=native

  svn:global-ignores
    *.diff
    *.patch

Properties on '/branches':
\end{verbatim}

\section{svnlook tree}\label{svn.ref.svnlook.c.tree}

\emph{Print the tree.}

\textbf{Synopsis}

\texttt{svnlook\ tree\ REPOS\_PATH\ {[}PATH\_IN\_REPOS{]}}

\subsection{Description}

Print the tree, starting at \textless PATH\_IN\_REPOS\textgreater{} (if
supplied; at the root of the tree otherwise), optionally showing node
revision IDs.

\subsection{Options}

\begin{verbatim}
--full-paths
--non-recursive (-N)
--revision (-r) REV
--show-ids
--transaction (-t) ID
\end{verbatim}

\subsection{Example}

This shows the tree output for revision 13 in our sample repository:

\begin{verbatim}
$ svnlook tree -r 13 /var/svn/repos
/
 trunk/
  button.c
  Makefile
  integer.c
 branches/
  bookstore/
   button.c
   Makefile
   integer.c
…
\end{verbatim}

Use the \texttt{-\/-show-ids} option to include node revision IDs
(unique internal identifiers for specific nodes in
Subversion\textquotesingle s versioned filesystem implementation):

\begin{verbatim}
$ svnlook tree -r 13 /var/svn/repos --show-ids
/ <0.0.r13/811>
 trunk/ <1.0.r9/551>
  button.c <2.0.r9/238>
  Makefile <3.0.r7/41>
  integer.c <4.0.r6/98>
 branches/ <5.0.r13/593>
  bookstore/ <1.1.r13/390>
   button.c <2.1.r12/85>
   Makefile <3.0.r7/41>
   integer.c <4.1.r13/109>
…
\end{verbatim}

For output which lends itself more readily to being parsed by scripts,
use the \texttt{-\/-full-paths} option, which causes \texttt{svnlook} to
print the full repository path of each tree item and to not use
indentation to indicate hierarchy:

\begin{verbatim}
$ svnlook tree -r 13 /var/svn/repos --show-ids --full-paths
/ <0.0.r13/811>
trunk/ <1.0.r9/551>
trunk/button.c <2.0.r9/238>
trunk/Makefile <3.0.r7/41>
trunk/integer.c <4.0.r6/98>
branches/ <5.0.r13/593>
branches/bookstore/ <1.1.r13/390>
branches/bookstore/button.c <2.1.r12/85>
branches/bookstore/Makefile <3.0.r7/41>
branches/bookstore/integer.c <4.1.r13/109>
…
\end{verbatim}

\section{svnlook uuid}\label{svn.ref.svnlook.c.uuid}

\emph{Print the repository\textquotesingle s UUID.}

\textbf{Synopsis}

\texttt{svnlook\ uuid\ REPOS\_PATH}

\subsection{Description}

Print the UUID for the repository. The UUID is the
repository\textquotesingle s \emph{u}niversal \emph{u}nique
\emph{id}entifier. The Subversion client uses this identifier to
differentiate between one repository and another.

\subsection{Options}

None

\subsection{Examples}

\begin{verbatim}
$ svnlook uuid /var/svn/repos
e7fe1b91-8cd5-0310-98dd-2f12e793c5e8
\end{verbatim}

\section{svnlook youngest}\label{svn.ref.svnlook.c.youngest}

\emph{Print the youngest revision number.}

\textbf{Synopsis}

\texttt{svnlook\ youngest\ REPOS\_PATH}

\subsection{Description}

Print the youngest revision number of a repository.

\subsection{Options}

None

\subsection{Examples}

This shows the youngest revision of our sample repository:

\begin{verbatim}
$ svnlook youngest /var/svn/repos/ 
42
\end{verbatim}

